\section{Two Execution Dependencies for Protection}\label{sec:evidence}
The field map connects two properties of Agent execution to defense design: information passes between services, and cover can become application work. Table~\ref{tab:conditional-risk} states the observation conditions under which particular signals matter.

\subsection{Feedback joins separately observable services}\label{sec:authorization}
The support-ticket execution exposes a selected tool destination, arguments received by that tool, a later model request and the model's response. WRA recovers query meaning from contacted domains under its name-association and supplied-trace conditions~\cite{jeong2026wra}. ANT, a network-traffic benchmark for Agent behavior, and the cloud benchmark measure task recognition from model traffic~\cite{li2026ant,pouryousef2026traffic}. The cloud benchmark reports ten-class balanced accuracy .853/.860 for graph/star coordination with all 247 features, and .540/.587 when size features are excluded and 100 of the 134 remaining features are selected. Both feature availability and selection change in that comparison. ANT's model-view values reconstructed from rounded figures remain approximate quantities in Supplement S2.

\begin{table*}[t]\centering\small
\caption{Conditional risk map. A row applies when its access and execution conditions hold. The final column states a protection priority or design hypothesis, not a demonstrated mitigation. Scores, targets and rights differ across rows, so there is no global danger ranking.}\label{tab:conditional-risk}
\setlength{\tabcolsep}{4pt}
\begin{tabularx}{\textwidth}{@{}P{.23\textwidth}P{.43\textwidth}Y@{}}\toprule
Interface and signal & Condition and evidenced inference & Protection consequence \\\midrule
\textbf{Tool egress:} selected domains and timing & The observer can associate contacted names with a supplied research trace. WRA recovers query meaning; linked histories use a separate persona endpoint~\cite{jeong2026wra}. & Prioritize destination association when the private target is research intent. Reduced name visibility needs a residual traffic assessment; executed decoys need task-utility constraints. \\\addlinespace[4pt]
\textbf{Model link:} lengths, directions and timing & Model traffic is visible and associated with the evaluated task. ANT and the cloud benchmark measure task information in this view~\cite{li2026ant,pouryousef2026traffic}. Initial instructions and later feedback coexist. & Tool-only protection leaves this view to be assessed. Compare model-only and joint adaptive attacks; isolate feedback before choosing a feedback-specific defense. \\\addlinespace[4pt]
\textbf{Coordination:} message graph, volume and sequence & Workflow records or coordination traffic expose execution organization. Composition and topology transfer tests change that organization under known task labels~\cite{dangol2026workflow,pouryousef2026traffic}. & Target observable scheduling and association, not just one payload size. Evaluate changed compositions and topology separately from recurrent-workflow privacy. \\\addlinespace[4pt]
\textbf{Client link:} packet counts, directions and bursts & AgentPrint's supplied browser interactions support hosted-agent identity inference~\cite{agentprint2025}. The result presumes those observation units and identities. & Client-facing traffic needs its own protection assessment even if tool links are hidden. Task discovery and new identities require separate tests. \\\addlinespace[4pt]
\textbf{Contacted recipient:} headers, commands or arguments & MARK, Trace and PrivacyTrace use recipient-visible representations. In PrivacyTrace, metadata alone is below its prior, while adding arguments improves attribute inference~\cite{kang2026whoseagent,ediga2026trace,duan2026privacytrace}. & Minimize sensitive request content and assess implementation signatures. Transport encryption does not conceal the recipient's intended plaintext. \\\addlinespace[4pt]
\textbf{Active probing:} challenges, choices and repairs & A party can elicit responses and observe later behavior. AgentTrace attributes models from repairs; AgentTell observes a choice under a disclosure rule~\cite{zhao2026agenttrace,shahriar2026agenttell}. & Test error normalization and restrictions on cross-context disclosures while retaining task recovery. These hypotheses need active-access evaluations. \\
\bottomrule\end{tabularx}\end{table*}


Recipient access changes the inference problem. PrivacyTrace's tool metadata/order yields attribute-inference ASR 14.53\%, below its 17.97\% schema-prior reference; adding arguments yields 27.71\%~\cite{duan2026privacytrace}. MARK's combined F1 .971 also uses browser behavior, while its protocol/header features are available at the contacted website~\cite{kang2026whoseagent}. Encryption against transit observation does not hide an addressed request from its recipient.

These results identify where task information is observable. They lead to a testable question: with the initial instruction fixed and task-bearing feedback released later, how much recognition remains on the model link under a stated tool-side policy? Section~\ref{sec:experiment} tests this question with constructed schemas and reports both the first-feedback and complete-execution views. Initial instructions, retained task fields and model responses are distinct possible contributors to the latter.

\paragraph{Defense consequence}
Enumerate the services and representations available to the attacker, then evaluate the views remaining after an intervention. A joint observer can use a single-link decision rule, even when a particular fitted joint classifier scores lower. At the application boundary, minimize forwarded context while preserving the information needed for the task; at the transport boundary, evaluate visible sizes, timing and association. Their effects belong to separate contrasts.

\subsection{Cover termination determines the execution constraint}\label{sec:crosslink}
Transport cover can terminate before business execution. Executed cover performs work and may affect resources, state or later decisions. The difference determines which utility must accompany a privacy measurement.

WRA's five research decoys reduce functional query recovery from .770 to .565~\cite[Sec.~8.1]{jeong2026wra}. Its utility scores 8.55 and 8.54 instead compare 100\% and 10\% destination visibility. Thus the utility observation does not describe the decoy intervention that produced the recovery reduction. WRA establishes additional research execution; it does not report harmful business-state writes. Matched protected-condition attacker examples also remain unresolved in the selected evidence.

The live OHTTP study refits a random forest on each protected capture and pairs six-class accuracy .30 with approximately thirteen-fold bytes and eleven-fold latency~\cite[Sec.~9.7]{dangol2026workflow}. Its separate nine-class .67 score belongs to another comparison. This design point binds adapted inference to transport cost; matched original-task utility was not located.

\paragraph{Task completion is only one execution outcome}
The requested refund can satisfy its target postcondition while another order, shared quota or ledger entry violates the application's state contract. Cover placement therefore determines what must be checked alongside completion. Section~\ref{sec:cover-study} illustrates this distinction with paired LLM executions and reports inference, state and resource outcomes together.

\paragraph{Defense consequence}
An executed-cover design needs a semantic constraint alongside a traffic budget. It can terminate decoys before business execution, or isolate their resources, state and planner feedback. These are testable design choices: measure adapted inference, additional work, original-task completion and collateral state changes on the same protected runs. The published WRA text supports the intervention/utility distinction; the OHTTP refit/cost example depends on one preprint.

\subsection{Execution changes delimit the two findings}\label{sec:targets}
A private recurring routine can be recognizable while its classifier transfers poorly to a new execution organization. In the communication-graph study, composition holdout keeps the nine workflow labels but changes their capability combinations: accuracy changes from about .67 to .18, with zero recall for seven classes~\cite[Sec.~9.12]{dangol2026workflow}. The cloud benchmark's graph/star change separately reports about .83 to .52 under its corresponding feature protocol~\cite{pouryousef2026traffic}. Each axis depends on one preprint.

Implementation and failure handling can also change the observation. Trace reports macro-F1 .981 in its mixed-session model-attribution test and .815 on scaffold holdout, retaining 1,010 of 1,050 planned clean runs~\cite{ediga2026trace}. AgentTrace's active probes reach 66.8\% attribution from initial outcomes and 96.0\% from ordered repairs at fifty probes~\cite{zhao2026agenttrace}. Excluding failed sessions and observing recoverable errors are different operations. These source-specific conditions delimit what an attack or defense says about new workloads; they do not estimate a common filtering effect.

Finally, a traffic target retains its own security meaning. ANT labels benchmark risk, AgentTell tests action choice under a disclosure rule, and CHeaT measures prevention of a defined CTF objective~\cite{li2026ant,shahriar2026agenttell,ayzenshteyn2025cheat}. Authorization additionally requires a principal and policy context. Two refunds that differ only in an unobserved approval can have identical traffic. This conditional argument identifies the extra policy input an authorization monitor needs.
